\section{Теоретическая часть}

\textbf{Разрывные колебания} --- это колебания, при которых медленные изменения состояния системы в некоторые моменты сменяются настолько быстрыми, что их можно рассматривать как мгновенные скачки. Они описываются уравнениями с малым параметром при производной, что позволяет применять метод разрывных колебаний.

\subsection{Модель системы}

Схема содержит туннельные диоды с $N$-образной вольт-амперной характеристикой $I(u)$, падающий участок которой $u_1 < u < u_2$ обеспечивает релаксационные колебания. Динамика схемы описывается уравнениями
\begin{align} \label{eq:th:1}
	\frac{C}{2}\dv{(u-E)}{t} + \frac{C}{2}\dv{u}{t} + i + I(u) = I(E-u),
\end{align}
\begin{align} \label{eq:th:2}
	L\dv{i}{t} = u - v, \qquad i = \frac{v-E}{2R}.
\end{align}

Переходя к безразмерному времени $\tau = Rt/L$, параметру $\mu = CR^2/L$ и переменным $x = 2u/E$, $y = 2iR/E$, получаем
\begin{align} \label{eq:th:3}
	\mu\dv{x}{\tau} = \varphi(x) - y, \qquad \dv{y}{\tau} = x - y,
\end{align}
где $\varphi(x) = \frac{2R}{E}\big[I(E/2 - u) + I(E/2 + u)\big]$. Вид функции $\varphi(x)$ зависит от напряжения питания $E$ и определяет режим работы схемы при $\mu \ll 1$.

\subsection{Быстрые и медленные движения}

При $\mu \ll 1$ движения на фазовой плоскости $(x, y)$ разделяются на медленные, происходящие в малой окрестности кривой $\varphi(x)$, и быстрые. Медленные движения описываются системой
\begin{align} \label{eq:th:4}
	y = \varphi(x), \qquad \dv{y}{\tau} = x - y,
\end{align}
а быстрые --- при $y = \text{const}$:
\begin{align} \label{eq:th:5}
	\dv{x}{\tau} = \frac{\varphi(x) - y}{\mu}.
\end{align}

На кривой медленных движений устойчивы только участки с $\varphi'(x) < 0$. Дойдя до границы участка, точка срывается и быстрым движением переходит на другой устойчивый участок кривой $\varphi(x)$. Чередование медленных и быстрых движений образует предельный цикл, то есть в системе устанавливаются автоколебания.

\subsection{Мультивибратор}

Мультивибратор --- генератор прямоугольных колебаний с короткими фронтами. Если на кривой $\varphi(x)$ остаётся один падающий участок, предельный цикл реализуется при любых начальных условиях (мягкий режим автоколебаний). Период и амплитуда определяются видом $\varphi(x)$ и параметром $\mu$.

\subsection{Триггер}

Триггер --- система с двумя устойчивыми состояниями равновесия, которая переводится из одного состояния в другое подачей запускающего импульса. Кривая $\varphi(x)$ имеет при этом два падающих участка, а области притяжения состояний разделены сепаратрисами.

При подаче импульса кривая $\varphi(x)$ смещается к виду, соответствующему режиму мультивибратора, и точка уходит из состояния равновесия. Слишком короткий импульс не успевает вывести её на участок медленных движений, а слишком длинный --- переводит дальше по кривой; в обоих случаях после снятия импульса система возвращается в исходное состояние. Поэтому существуют минимальная $\tau_{\min}$ и максимальная $\tau_{\max}$ допустимые длительности запускающего импульса, а также пороговые значения его амплитуды.

Разрешающее время триггера --- наименьший интервал между импульсами, при котором переключение устойчиво; если период следования меньше него, схема переключается через импульс, то есть делит частоту.

\subsection{Кипп-реле}

Кипп-реле (одновибратор) --- спусковая схема с одним устойчивым состоянием, которая на короткий запускающий импульс отвечает широким прямоугольным импульсом, длительность которого определяется только параметрами схемы. Её динамика описывается уравнениями
\begin{align} \label{eq:th:6}
	L\dv{i}{t} + Ri = E - u, \qquad I(u) + C\dv{u}{t} = i + i_\text{имп}.
\end{align}

В безразмерных переменных $x = u/E$, $y = Ri/E$, $e = i_\text{имп}R/E$ и при $\mu = CR^2/L \ll 1$ получаем
\begin{align} \label{eq:th:7}
	\mu\dv{x}{\tau} = y - \varphi(x) + e, \qquad \dv{y}{\tau} = 1 - x - y,
\end{align}
а медленные движения описываются системой
\begin{align} \label{eq:th:8}
	y = \varphi(x) + e, \qquad \dv{y}{\tau} = 1 - x - y.
\end{align}

Без входного сигнала ($e = 0$) схема находится в устойчивом равновесии. При подаче импульса ($e > 0$) кривая медленных движений смещается, а изображающая точка после быстрого скачка медленно движется к новому положению равновесия. Если импульс закончится раньше, схема вернётся в исходное состояние, поэтому длительность выходного импульса определяется движением по кривой медленных движений и не зависит от длительности запускающего.
